\providecommand{\Res}[2]{\qty[round-mode=figures,round-precision=3]{\fpeval{#1}}{#2}}
\providecommand{\Val}[1]{\num[round-mode=figures,round-precision=4]{\fpeval{#1}}}
\providecommand{\ResSci}[2]{\qty[round-mode=figures,round-precision=3,exponent-mode=scientific]{\fpeval{#1}}{#2}}
\subsubsection*{Question a) Longueur d'onde}
\begin{enumerate}[label=\arabic*)]
    \item \textbf{Variables connues et recherchées :}
    \begin{itemize}[label= -]
        \item $f = \qty{\v{f}}{Hz}$, $v = \qty{\v{v}}{m/s}$
        \item Variable recherchée : $\lambda$ (m)
    \end{itemize}
    \item \textbf{Équations pertinentes :}
    \[ v = \lambda f \]
    \item \textbf{Isolation de la variable recherchée :}
    \[ \lambda = \frac{v}{f} \]
    \item \textbf{Résolution numérique :}
    \[ \lambda = \frac{\num{\v{v}}}{\num{\v{f}}} = \Res{\v{v}/\v{f}}{m} \]
    \item \textbf{Réponse :} la longueur d'onde est de \textbf{\Res{\v{v}/\v{f}}{m}}.
\end{enumerate}

\subsubsection*{Question b) Déphasage entre deux points}
\begin{enumerate}[label=\arabic*)]
    \item \textbf{Variables connues et recherchées :}
    \begin{itemize}[label= -]
        \item $\Delta x = \qty{\vsi{dx}}{m}$, $\lambda = \Res{\v{v}/\v{f}}{m}$
        \item Variable recherchée : $\Delta\phi$ (rad)
    \end{itemize}
    \item \textbf{Équations pertinentes :}
    \[ \Delta\phi = k\,\Delta x = \frac{2\pi}{\lambda}\,\Delta x \]
    \item \textbf{Isolation de la variable recherchée :} l'équation est déjà isolée.
    \item \textbf{Résolution numérique :}
    \[ \Delta\phi = \frac{2\pi}{\Val{\v{v}/\v{f}}} \cdot \num{\vsi{dx}} = \Res{2*pi*\v{f}/\v{v}*\vsi{dx}}{rad} \]
    \item \textbf{Réponse :} le déphasage est de \textbf{\Res{2*pi*\v{f}/\v{v}*\vsi{dx}}{rad}}.
\end{enumerate}

\subsubsection*{Question c) Déphasage entre deux instants}
\begin{enumerate}[label=\arabic*)]
    \item \textbf{Variables connues et recherchées :}
    \begin{itemize}[label= -]
        \item $\Delta t = \qty{\vsi{dt}}{s}$, $f = \qty{\v{f}}{Hz}$
        \item Variable recherchée : $\Delta\phi$ (rad)
    \end{itemize}
    \item \textbf{Équations pertinentes :}
    \[ \Delta\phi = \omega\,\Delta t = 2\pi f\,\Delta t \]
    \item \textbf{Isolation de la variable recherchée :} l'équation est déjà isolée.
    \item \textbf{Résolution numérique :}
    \[ \Delta\phi = 2\pi \cdot \num{\v{f}} \cdot \num{\vsi{dt}} = \Res{2*pi*\v{f}*\vsi{dt}}{rad} \]
    \item \textbf{Réponse :} le déphasage est de \textbf{\Res{2*pi*\v{f}*\vsi{dt}}{rad}}.
\end{enumerate}

\subsubsection*{Question d) Opposition de phase}
\begin{enumerate}[label=\arabic*)]
    \item \textbf{Variables connues et recherchées :} $\lambda = \Res{\v{v}/\v{f}}{m}$ ; variable recherchée : $\Delta x_{\min}$
    \item \textbf{Équations pertinentes :} opposition de phase : $\Delta\phi = \pi$, donc $\dfrac{2\pi}{\lambda}\Delta x = \pi$.
    \item \textbf{Isolation de la variable recherchée :}
    \[ \Delta x_{\min} = \frac{\lambda}{2} \]
    \item \textbf{Résolution numérique :}
    \[ \Delta x_{\min} = \frac{\Val{\v{v}/\v{f}}}{2} = \Res{\v{v}/(2*\v{f})}{m} \]
    \item \textbf{Réponse :} la distance minimale est de \textbf{\Res{100*\v{v}/(2*\v{f})}{cm}}.
\end{enumerate}
